\documentclass{scrartcl}
\usepackage{graphicx}  % required for inserting images

% for cyrillic and Bulgarian language
\usepackage[utf8]{inputenc}
\usepackage[T2A]{fontenc}
\usepackage[bulgarian]{babel}

% math packages
\usepackage{amsmath}
\usepackage{float}
\usepackage{amsfonts}
\usepackage{amssymb}
\usepackage{amsthm}
\usepackage{cancel}

\usepackage{ragged2e}  % adds FlushLeft

\usepackage{listings}
\usepackage{xcolor}

\definecolor{codegreen}{rgb}{0,0.6,0}
\definecolor{codegray}{rgb}{0.5,0.5,0.5}
\definecolor{codepurple}{rgb}{0.58,0,0.82}
\definecolor{backcolour}{rgb}{0.95,0.95,0.92}

\lstdefinestyle{mystyle}{
    language=Octave,
    backgroundcolor=\color{backcolour},   
    commentstyle=\color{codegreen},
    keywordstyle=\color{magenta},
    numberstyle=\tiny\color{codegray},
    stringstyle=\color{codepurple},
    basicstyle=\ttfamily\footnotesize,
    breakatwhitespace=false,         
    breaklines=true,
    captionpos=b,
    keepspaces=false,                 
    numbersep=5pt,
    showspaces=false,
    showstringspaces=false,
    showtabs=false,
    tabsize=2,
    columns=fullflexible,
}

\lstset{style=mystyle}

\usepackage[colorlinks=true, linkcolor=blue, urlcolor=cyan]{hyperref}
\usepackage{tikz}
\usetikzlibrary{arrows.meta}

\newtheorem{theorem}{Теорема}
\newtheorem{definition}{Дефиниция}

\newtheorem{lemma}{Лема}
\newtheorem{example}{Пример}
\newtheorem{remark}{Забележка}

\title{Теорема на Ферма. Теорема за средни стойности. Формула на Тейлър}
\author{Ивайло Андреев}
\date{1 юли 2026}

\begin{document}
\maketitle

\tableofcontents

\newpage

\begin{FlushLeft}

\section{Локални екстремуми и Теорема на Ферма}

\begin{definition}[Локален екстремум]
Нека $f: D \to \mathbb{R}$ е функция, дефинирана върху множеството $D \subset \mathbb{R}$, и $x_0 \in D$.
\begin{enumerate}
    \item Казваме, че $f$ има \textbf{локален минимум} в точката $x_0$, ако съществува $\delta > 0$ такова, че околността $(x_0 - \delta, x_0 + \delta) \subset D$ и за всяко $x \in (x_0 - \delta, x_0 + \delta)$ е изпълнено:
    \[ f(x_0) \le f(x) \]
    \item Казваме, че $f$ има \textbf{локален максимум} в точката $x_0$, ако съществува $\delta > 0$ такова, че околността $(x_0 - \delta, x_0 + \delta) \subset D$ и за всяко $x \in (x_0 - \delta, x_0 + \delta)$ е изпълнено:
    \[ f(x_0) \ge f(x) \]
\end{enumerate}
Локалните минимуми и максимуми се обединяват под общото наименование \textbf{локални екстремуми}.
\end{definition}

\begin{theorem}[Теорема на Ферма - Необходимо условие за локален екстремум]
Нека $f: D \to \mathbb{R}$ е функция, дефинирана върху множеството $D \subset \mathbb{R}$, и $x_0 \in D$ и нека $f$ има локален екстремум в $x_0$. Ако $f$ е диференцируема в точката $x_0$, то:
\[ f'(x_0) = 0 \]
\end{theorem}

\begin{proof}
Ще проведем доказателството за случая, когато $f$ има \textbf{локален максимум} в точката $x_0$ (случаят за локален минимум се доказва напълно аналогично).

По дефиниция, щом $f$ има локален максимум в $x_0 \in D$, то съществува околност $(x_0 - \delta, x_0 + \delta)$, в която:
\[ f(x) \le f(x_0) \quad \iff \quad f(x) - f(x_0) \le 0 \]
за всяко $x$ от тази околност.

Разглеждаме производната на $f$ в точката $x_0$:

\[ f'(x_0) = \lim_{h \to 0} \frac{f(x_0 + h) - f(x_0)}{h} \]

Тъй като тази граница съществува, левите и десните едностранни граници също съществуват и са равни на $f'(x_0)$. Нека изследваме знака на диференчното частно за $h \in (-\delta, \delta) \setminus \{0\}$:

\begin{enumerate}
    \item \textbf{Случай 1 (Дясна граница):} Нека $h \in (0, \delta)$. Тогава:
    \[ \frac{f(x_0 + h) - f(x_0)}{h} = \frac{\le 0}{> 0} \le 0 \]
    Преминавайки към граница при $h \to +0$, от теоремата за запазване на знака при граничен преход получаваме:
    \[ f'(x_0) = \lim_{h \to +0} \frac{f(x_0 + h) - f(x_0)}{h} \le 0 \]
    
    \item \textbf{Случай 2 (Лява граница):} Нека $h \in (-\delta, 0)$. Тогава:
    \[ \frac{f(x_0 + h) - f(x_0)}{h} = \frac{\le 0}{< 0} \ge 0 \]
    Преминавайки към граница при $h \to -0$, получаваме:
    \[ f'(x_0) = \lim_{h \to -0} \frac{f(x_0 + h) - f(x_0)}{h} \ge 0 \]
\end{enumerate}

Получаваме $f'(x_0) \le 0$ и $f'(x_0) \ge 0$ следва

\[ f'(x_0) = 0 \]
\end{proof}

\begin{remark}[Геометричен смисъл]
Геометрично теоремата на Ферма означава, че в точката на локален екстремум допирателната към графиката на функцията $y = f(x)$ е хоризонтална права, т.е. успоредна на абсцисната ос $Ox$.
\end{remark}

\newpage

\section{Теореми за средните стойности}

\begin{theorem}[Теорема на Вайерщрас]
Нека $f: [a, b] \to \mathbb{R}$ е непрекъсната. Тогава $f$ достига в $[a, b]$ своята най-голяма стойност (максимум) $M$ и своята най-малка (минимум) $m$.

Тоест, съществуват точки $x_1, x_2 \in [a, b]$, такива че $f(x_1) = m$ и $f(x_2) = M$ и $m \le f(x) \le M \; \forall x \in [a, b]$.
\end{theorem}

\subsection{Теорема на Рол}

\begin{theorem}[Теорема на Рол]
Нека функцията $f(x)$ отговаря на следните три условия:
\begin{enumerate}
    \item $f(x)$ е непрекъсната в затворения интервал $[a, b]$;
    \item $f(x)$ притежава производна (диференцируема е) поне в отворения интервал $(a, b)$;
    \item Стойностите в краищата на интервала са равни: $f(a) = f(b)$.
\end{enumerate}
Тогава съществува поне една точка $c \in (a, b)$, такава че:
\[ f'(c) = 0 \]
\end{theorem}

\begin{proof}
Геометричният смисъл е показан по-долу за двата случая в доказателството:

\begin{figure}[H]
\centering
\begin{tikzpicture}[scale=1.1]
  % Case 1: constant function
  \begin{scope}[xshift=0cm]
    \draw[-{Stealth}] (-0.3,0) -- (4,0) node[right] {$x$};
    \draw[-{Stealth}] (0,-0.3) -- (0,2.5) node[above] {$y$};
    \draw[blue, thick] (0.5,1.5) -- (3.5,1.5);
    \filldraw (0.5,1.5) circle (1.5pt);
    \filldraw (3.5,1.5) circle (1.5pt);
    \draw[dashed] (0.5,0) node[below] {$a$} -- (0.5,1.5);
    \draw[dashed] (3.5,0) node[below] {$b$} -- (3.5,1.5);
    \node[left] at (0,1.5) {$c$};
    \draw[dashed] (0,1.5) -- (0.5,1.5);
    \node[below, align=center] at (2,-0.6) {Случай 1: $f = \text{const}$,\\ $f'(x) = 0$ навсякъде};
  \end{scope}
  % Case 2: function with visible extremum
  \begin{scope}[xshift=5.5cm]
    \draw[-{Stealth}] (-0.3,0) -- (4,0) node[right] {$x$};
    \draw[-{Stealth}] (0,-0.3) -- (0,2.5) node[above] {$y$};
    \draw[blue, thick] plot[domain=0.5:3.5, samples=60]
      (\x, {0.9 + 1.4*sin(deg(3.14159*( \x - 0.5)/3))});
    \filldraw (0.5,0.9) circle (1.5pt);
    \filldraw (3.5,0.9) circle (1.5pt);
    \draw[dashed] (0.5,0) node[below] {$a$} -- (0.5,0.9);
    \draw[dashed] (3.5,0) node[below] {$b$} -- (3.5,0.9);
    % maximum point
    \draw[dashed] (2,0) node[below] {$c$} -- (2,2.3);
    \filldraw[red] (2,2.3) circle (1.5pt) node[above] {$M$};
    % horizontal tangent
    \draw[red, thick] (1.2,2.3) -- (2.8,2.3) node[right] {$f'(c)=0$};
    \node[below, align=center] at (2,-0.6) {Случай 2: $m < M$,\\ екстремумът е вътрешен};
  \end{scope}
\end{tikzpicture}
\caption{Геометричен смисъл на теоремата на Рол — Случай 1: $f = \mathrm{const}$, $f' \equiv 0$; Случай 2: екстремумът се достига във вътрешна точка $c$, откъдето $f'(c) = 0$}
\end{figure}

Тъй като $f(x)$ е непрекъсната в крайния и затворен интервал $[a, b]$, съгласно Теоремата на Вайерщрас тя достига своя минимум $m$ и своя максимум $M$ в този интервал. Нека:
\[ m = \min_{x \in [a, b]} f(x) = f(x_1), \quad M = \max_{x \in [a, b]} f(x) = f(x_2) \quad (x_1, x_2 \in [a, b]) \]

Разглеждаме два основни случая за стойностите на $m$ и $M$:

\begin{enumerate}
    \item \textbf{Случай 1:} $m = M$.
    Ако минимумът е равен на максимума, то функцията е константа в целия интервал: $f(x) = c = \text{const}$ за всяко $x \in [a, b]$. Тогава нейната производна е нула във всяка точка от интервала: $f'(x) = 0$. В този случай всяка точка от отворения интервал $(a, b)$ може да бъде избрана за точка $c$.
    
    \item \textbf{Случай 2:} $m < M$.
    Тъй като по условие $f(a) = f(b)$, не е възможно и двете стойности $m$ и $M$ да се достигат едновременно в краищата на интервала. Наистина, ако и $m$, и $M$ се достигаха само в краищата, то от $f(a) = f(b)$ би следвало $m = M$, което противоречи на предположението $m < M$. Следователно, поне една от двете стойности (или минимумът $m$, или максимумът $M$) се достига във вътрешна точка $c \in (a, b)$.
    
    Тъй като $c$ е вътрешна точка на интервала, тя се явява точка на локален екстремум (локален минимум или локален максимум) за функцията $f$. Освен това, по условие функцията е диференцируема в отворения интервал $(a, b)$, следователно тя е диференцируема и в точката $c$. 
    
    Прилагаме директно \textbf{Теоремата на Ферма} за точката $c$, откъдето следва:
    \[ f'(c) = 0 \]
\end{enumerate}
С това теоремата на Рол е доказана.
\end{proof}

\subsection{Теорема на Лагранж}

\begin{theorem}[Теорема на Лагранж]
Нека функцията $f(x)$ е непрекъсната в затворения интервал $[a, b]$ и диференцируема в отворения интервал $(a, b)$. Тогава съществува поне една точка $c \in (a, b)$, такава че:
\[ f(b) - f(a) = f'(c)(b - a) \]
\end{theorem}

\begin{proof}
Геометричният смисъл е показан по-долу — теоремата гарантира \textbf{поне една} вътрешна точка $c$, в която наклонът на допирателната съвпада с наклона на секущата през $(a, f(a))$ и $(b, f(b))$ (в примера по-долу се виждат две такива точки $c_1$ и $c_2$):

\begin{figure}[H]
\centering
\begin{tikzpicture}[scale=1.2]
  \draw[-{Stealth}] (-0.3,0) -- (5.5,0) node[right] {$x$};
  \draw[-{Stealth}] (0,-0.3) -- (0,3.2) node[above] {$y$};
  % curve: f(x) = 0.18*(x-0.6)*(x-4.4)*(x-2.5) + 0.5*(x-0.6) + 0.8
  % f(0.6) = 0.8,  f(4.4) = 2.7
  \draw[blue, thick] plot[domain=0.6:4.4, samples=80]
    (\x, {0.18*(\x-0.6)*(\x-4.4)*(\x-2.5) + 0.5*(\x-0.6) + 0.8});
  % endpoints A=(0.6, 0.8)  B=(4.4, 2.7)
  \coordinate (A) at (0.6, 0.8);
  \coordinate (B) at (4.4, 2.7);
  \filldraw (A) circle (1.5pt);
  \filldraw (B) circle (1.5pt);
  \draw[dashed] (0.6,0) node[below] {$a$} -- (A);
  \draw[dashed] (4.4,0) node[below] {$b$} -- (B);
  \draw[dashed] (0,0.8) node[left] {$f(a)$} -- (A);
  \draw[dashed] (0,2.7) node[left] {$f(b)$} -- (B);
  % secant: y = 0.5x + 0.5  (passes exactly through A and B)
  \draw[gray, thick] (0.0, 0.5) -- (5.0, 3.0) node[right] {секуща};

  % f'(x) = slope of secant = 0.5  at  c1≈1.4  and  c2≈3.6
  % f(1.4)  = 0.18*(0.8)*(-3)*(-1.1) + 0.5*0.8 + 0.8 = 0.4752+1.2 = 1.6752
  % f(3.6)  = 0.18*(3)*(-0.8)*(1.1) + 0.5*3 + 0.8   = -0.4752+2.3 = 1.8248
  % tangent at c1=(1.4, 1.6752): y = 0.5x + 0.9752
  \filldraw[red] (1.4, 1.6752) circle (1.5pt);
  \draw[dashed] (1.4, 0) node[below] {$c_1$} -- (1.4, 1.6752);
  \draw[red, thick] (0.5, 1.2252) -- (2.3, 2.1252);
  % tangent at c2=(3.6, 1.8248): y = 0.5x + 0.0248
  \filldraw[red] (3.6, 1.8248) circle (1.5pt);
  \draw[dashed] (3.6, 0) node[below] {$c_2$} -- (3.6, 1.8248);
  \draw[red, thick] (2.7, 1.375) -- (4.5, 2.275)
    node[right] {$f'(c) = \dfrac{f(b)-f(a)}{b-a}$};
\end{tikzpicture}
\caption{Геометричен смисъл на теоремата на Лагранж — теоремата гарантира поне една точка $c$; в примера са видими две: $c_1$ и $c_2$}
\end{figure}

За доказателството ще използваме теоремата на Рол за подходящо избрана функция.

Нека дефинираме числото $k$, представляващо наклона на секущата през краищата на графиката $(a, f(a))$ и $(b, f(b))$:

\[ k = \frac{f(b) - f(a)}{b - a} \]

Конструираме помощната функция $F(x)$ по следния начин:

\[ F(x) = f(x) - k \cdot x = f(x) - \frac{f(b) - f(a)}{b - a} \cdot x \]

Проверяваме условията на теоремата на Рол за функцията $F(x)$ в интервала $[a, b]$:
\begin{enumerate}
    \item $F(x)$ е непрекъсната в $[a, b]$, защото е разлика на непрекъснатата функция $f(x)$ и линейната функция $k \cdot x$.
    \item $F(x)$ е диференцируема в $(a, b)$, тъй като е разлика на диференцируеми функции. Нейната производна е:
    \[ F'(x) = f'(x) - k = f'(x) - \frac{f(b) - f(a)}{b - a} \]
    \item Проверяваме стойностите в краищата на интервала:
    \[ F(a) = f(a) - \frac{f(b) - f(a)}{b - a} \cdot a = \frac{f(a)b - f(a)a - f(b)a + f(a)a}{b - a} = \frac{f(a)b - f(b)a}{b - a} \]
    \[ F(b) = f(b) - \frac{f(b) - f(a)}{b - a} \cdot b = \frac{f(b)b - f(a)b - f(b)b + f(a)b}{b - a} = \frac{f(a)b - f(b)a}{b - a} \]
    Виждаме, че $F(a) = F(b)$.
\end{enumerate}

Тъй като всички условия на теоремата на Рол са изпълнени за $F(x)$, то съществува точка $c \in (a, b)$, такава че $F'(c) = 0$. Тоест:

\[ F'(c) = f'(c) - \frac{f(b) - f(a)}{b - a} = 0 \quad \Rightarrow \quad f'(c) = \frac{f(b) - f(a)}{b - a} \]

Умножавайки двете страни на равенството по $(b - a)$, получаваме:

\[ f(b) - f(a) = f'(c)(b - a) \]
\end{proof}

\subsection{Теорема на Коши}

\begin{theorem}[Теорема на Коши]
Нека функциите $f(x)$ и $g(x)$ са непрекъснати в затворения интервал $[a, b]$ и диференцируеми в отворения интервал $(a, b)$, като $g'(x) \neq 0$ за всяко $x \in (a, b)$. Тогава съществува точка $c \in (a, b)$, такава че:
\[ \frac{f(b) - f(a)}{g(b) - g(a)} = \frac{f'(c)}{g'(c)} \]
\end{theorem}

\begin{proof}
[Геометричният смисъл е опционален] Геометричният смисъл на теоремата на Коши се разкрива чрез параметричната крива $(g(t), f(t))$ във \textbf{фазовата равнина}: нека $A = (g(a), f(a))$ и $B = (g(b), f(b))$ са краищата на кривата, а $C = (g(c), f(c))$ е вътрешна точка, в която допирателната е паралелна на секущата $AB$.

\begin{figure}[H]
\centering
\begin{tikzpicture}[scale=1.5]
  \draw[-{Stealth}] (-0.3,0) -- (2.8,0) node[right] {$g(x)$};
  \draw[-{Stealth}] (0,-0.3) -- (0,1.8) node[above] {$f(x)$};
  % parametric curve (2cos(t), sin(t)) -- quarter ellipse
  \draw[blue, thick] plot[domain=0:90, samples=60]
    ({2*cos(\x)}, {sin(\x)});
  \draw[-{Stealth}, blue, thick] ({2*cos(48)}, {sin(48)}) -- ({2*cos(50)}, {sin(50)});
  % endpoints
  \filldraw (2,0) circle (1.2pt) node[below] {$A$};
  \filldraw (0,1) circle (1.2pt) node[left] {$B$};
  % secant through A=(2,0) and B=(0,1): y = -x/2 + 1
  \draw[gray, thick] (-0.1, 1.05) -- (2.3, -0.15);
  % tangent point C = (sqrt(2), sqrt(2)/2), tangent: y = -x/2 + sqrt(2)
  \filldraw[red] ({sqrt(2)}, {sqrt(2)/2}) circle (1.2pt) node[above=2pt] {$C$};
  \draw[dashed] ({sqrt(2)}, 0) -- ({sqrt(2)}, {sqrt(2)/2});
  \draw[dashed] (0, {sqrt(2)/2}) -- ({sqrt(2)}, {sqrt(2)/2});
  \draw[red, thick] (0.3, {-0.15 + sqrt(2)}) -- (2.2, {-1.1 + sqrt(2)});
\end{tikzpicture}
\caption{Геометричен смисъл на теоремата на Коши — в точката $C$ допирателната е паралелна на секущата $AB$, откъдето $\frac{f'(c)}{g'(c)} = \frac{f(b)-f(a)}{g(b)-g(a)}$}
\end{figure}

Доказателството се състои от две стъпки:

\textbf{Стъпка 1: Доказателство, че знаменателят $g(b) - g(a) \neq 0$.}
Ще докажем това твърдение с допускане на противното. Да допуснем, че $g(b) - g(a) = 0$, което означава $g(a) = g(b)$.
За функцията $g(x)$ са изпълнени всички условия на теоремата на Рол в интервала $[a, b]$ (тя е непрекъсната в $[a, b]$, диференцируема в $(a, b)$ и $g(a) = g(b)$). Тогава, по теоремата на Рол, трябва да съществува точка $\xi \in (a, b)$, за която $g'(\xi) = 0$.
Това обаче директно противоречи на условието на теоремата, че $g'(x) \neq 0$ за всяко $x \in (a, b)$. Следователно допускането е невярно и:
\[ g(b) - g(a) \neq 0 \]

\textbf{Стъпка 2: Използване на помощна функция.}
Дефинираме помощната функция $H(x)$:

\[ H(x) = f(x) - \lambda \cdot g(x), \quad \text{където } \lambda = \frac{f(b) - f(a)}{g(b) - g(a)} \]

Проверяваме условията на Рол за $H(x)$ в $[a, b]$:
\begin{enumerate}
    \item $H(x)$ е непрекъсната в $[a, b]$ и диференцируема в $(a, b)$ като линейна комбинация от непрекъснати и диференцируеми функции. Нейната производна е $H'(x) = f'(x) - \lambda \cdot g'(x)$.
    \item Проверяваме равенството в краищата:
    \[ H(a) = f(a) - \frac{f(b) - f(a)}{g(b) - g(a)} \cdot g(a) = \frac{f(a)g(b) - f(a)g(a) - f(b)g(a) + f(a)g(a)}{g(b) - g(a)} = \frac{f(a)g(b) - f(b)g(a)}{g(b) - g(a)} \]
    \[ H(b) = f(b) - \frac{f(b) - f(a)}{g(b) - g(a)} \cdot g(b) = \frac{f(b)g(b) - f(b)g(a) - f(b)g(b) + f(a)g(b)}{g(b) - g(a)} = \frac{f(a)g(b) - f(b)g(a)}{g(b) - g(a)} \]
    Следователно $H(a) = H(b)$.
\end{enumerate}

Прилагаме теоремата на Рол за $H(x)$ в интервала $[a, b]$, откъдето следва, че съществува точка $c \in (a, b)$, за която $H'(c) = 0$:
\[ H'(c) = f'(c) - \lambda \cdot g'(c) = 0 \quad \Rightarrow \quad f'(c) = \frac{f(b) - f(a)}{g(b) - g(a)} \cdot g'(c) \]
Понеже $c \in (a, b)$, от условието $g'(x) \neq 0$ за всяко $x \in (a, b)$ следва $g'(c) \neq 0$. Следователно можем да разделим двете страни на $g'(c)$ и получаваме:

\[ \frac{f(b) - f(a)}{g(b) - g(a)} = \frac{f'(c)}{g'(c)} \]
С това теоремата е напълно доказана.
\end{proof}

\newpage

\section{Формула на Тейлър с остатъчен член във формата на Лагранж}

Формулата на Тейлър ни позволява да апроксимираме достатъчно гладки функции чрез полиноми.

\begin{definition}[Полином на Тейлър]
Нека функцията $f$ е $n-1$ пъти диференцируема в околност на точката $x_0$ и $n$ пъти диференцируема в точката $x_0$. Полиномът:

\[ T_n(x) = f(x_0) + \frac{f'(x_0)}{1!}(x - x_0) + \frac{f''(x_0)}{2!}(x - x_0)^2 + \dots + \frac{f^{(n)}(x_0)}{n!}(x - x_0)^n = \sum_{k=0}^{n} \frac{f^{(k)}(x_0)}{k!}(x - x_0)^k \]

се нарича \textbf{Полином на Тейлър} от степен $n$ за функцията $f$ около точката $x_0$.
\end{definition}

Ако означим разликата между функцията и нейния полином на Тейлър с $R_n(x)$, получаваме \textbf{Формулата на Тейлър}:

\[ f(x) = T_n(x) + R_n(x) \]

където $R_n(x) = f(x) - T_n(x)$ се нарича \textbf{остатъчен член}. За да бъде формулата полезна, ни е необходима точна оценка за този остатъчен член.

\begin{theorem}[Формула на Тейлър с остатъчен член на Лагранж]
Нека функцията $f$ е $n+1$ пъти диференцируема в отворен интервал, съдържащ точката $x_0$. Нека $x$ е от интервала. Тогава съществува точка $c$ строго между $x_0$ и $x$, такава че остатъчният член $R_n(x)$ има вида:

\[ R_n(x) = \frac{f^{(n+1)}(c)}{(n+1)!}(x - x_0)^{n+1} \]
\end{theorem}

\begin{proof}
\textbf{Скица на доказателството. [За яснота на читателя, не е нужно за темата]}
\begin{enumerate}
    \item Въвеждаме две подходящо избрани помощни функции $\varphi(t)$ и $\psi(t)$.
    \item Изследваме техните производни и опростяваме получените изрази.
    \item Определяме стойностите им на краищата на разглеждания интервал.
    \item Проверяваме, че е приложима Теоремата на Коши.
    \item Прилагаме Теоремата на Коши, за да свържем стойностите и производните им в вътрешна точка.
    \item От полученото извеждаме обща форма за остатъчния член и чрез специален избор на параметър получаваме формата на Лагранж.
\end{enumerate}

Фиксираме точката $x \neq x_0$. Разглеждаме интервала $I$ с краища $x_0$ и $x$. Дефинираме помощна функция $\varphi(t)$ за $t \in I$ по следния начин:
\[ \varphi(t) = f(x) - \sum_{k=0}^{n} \frac{f^{(k)}(t)}{k!}(x - t)^k \]
Разписваме подробно сумата:
\[ \varphi(t) = f(x) - \left[ f(t) + \frac{f'(t)}{1!}(x - t) + \frac{f''(t)}{2!}(x - t)^2 + \dots + \frac{f^{(n)}(t)}{n!}(x - t)^n \right] \]

Нека пресметнем стойностите на функцията $\varphi(t)$ в краищата на интервала ($t = x_0$ и $t = x$):
\begin{enumerate}
    \item За $t = x$: Всички членове от вида $(x - t)^k$ стават нула за $k \ge 1$. Остава само:
    \[ \varphi(x) = f(x) - f(x) = 0 \]
    \item За $t = x_0$: Изразът съвпада точно с дефиницията за остатъчния член $R_n(x)$:
    \[ \varphi(x_0) = f(x) - T_n(x) = R_n(x) \]
\end{enumerate}

Сега диференцираме $\varphi(t)$:

\[ \varphi'(t) = 0 - \left[ f'(t) + \left( \frac{f''(t)}{1!}(x-t) - f'(t) \right) + \left( \frac{f^{(3)}(t)}{2!}(x-t)^2 - \frac{f''(t)}{1!}(x-t) \right) + \dots \right] \]

Забелязваме, че се получава телескопична сума, в която членовете (без последния) се унищожават по двойки. Така получаваме:

\[ \varphi'(t) = -\frac{f^{(n+1)}(t)}{n!}(x - t)^n \]

Въвеждаме втора помощна функция $\psi(t) = (x - t)^p$, където $p \ge 1, p \in \mathbb{N}$ е параметър. Производната на $\psi(t)$ е:

\[ \psi'(t) = -p(x - t)^{p-1} \]

Понеже за всяко $t$ от вътрешността на интервала е в сила, че $t \neq x$, следва че $\psi'(t) = -p(x-t)^{p-1} \neq 0$.

За функциите $\varphi(t)$ и $\psi(t)$ са изпълнени условията на \textbf{Теоремата на Коши} в интервала с краища $x_0$ и $x$. Следователно съществува вътрешна точка $c$ между $x_0$ и $x$, такава че:

\[ \frac{\varphi(x) - \varphi(x_0)}{\psi(x) - \psi(x_0)} = \frac{\varphi'(c)}{\psi'(c)} \]

Заместваме известните стойности $\varphi(x) = 0$, $\varphi(x_0) = R_n(x)$, $\psi(x) = 0$ и $\psi(x_0) = (x - x_0)^p$:

\[ \frac{0 - R_n(x)}{0 - (x - x_0)^p} = \frac{-\frac{f^{(n+1)}(c)}{n!}(x - c)^n}{-p(x - c)^{p-1}} \]
\[ \frac{R_n(x)}{(x - x_0)^p} = \frac{f^{(n+1)}(c)}{n! \cdot p} (x - c)^{n - p + 1} \]

Оттук изразяваме общата форма на остатъчния член:

\[ R_n(x) = \frac{f^{(n+1)}(c)}{n! \cdot p} (x - x_0)^p (x - c)^{n - p + 1} \]

Избираме стойност на параметъра $p = n + 1$:

\[ R_n(x) = \frac{f^{(n+1)}(c)}{n! \cdot (n + 1)} (x - x_0)^{n+1} (x - c)^{n - (n+1) + 1} \]

Тъй като $n! \cdot (n+1) = (n+1)!$ и $(x - c)^0 = 1$, изразът се опростява до:

\[ R_n(x) = \frac{f^{(n+1)}(c)}{(n+1)!}(x - x_0)^{n+1} \]
\end{proof}

[Примерът не е нужен за темата]
\begin{example}[Развиване на $\sin x$ в ред на Тейлър около $x_0 = 0$]
Пресмятаме производните на $f(x) = \sin x$ в $x_0 = 0$:

\[
f(0)=0,\quad f'(0)=1,\quad f''(0)=0,\quad f'''(0)=-1,\quad f^{(4)}(0)=0,\quad \ldots
\]

Производните се повтарят с период 4. Тъй като $\sin x$ е \textbf{нечетна функция}, всички четни производни в $0$ са нула, т.е. в реда участват само нечетни степени. Заместваме в полинома на Тейлър:

\[
\sin x = \sum_{k=0}^{\infty} \frac{(-1)^k}{(2k+1)!}\, x^{2k+1}
       = x - \frac{x^3}{3!} + \frac{x^5}{5!} - \frac{x^7}{7!} + \cdots
\]

Първите три последователни приближения са:

\begin{align*}
T_1(x) &= x \\[3pt]
T_3(x) &= x - \dfrac{x^3}{6} \\[3pt]
T_5(x) &= x - \dfrac{x^3}{6} + \dfrac{x^5}{120}
\end{align*}

Всяко следващо приближение „следи" $\sin x$ в по-широка околност на $x_0 = 0$:

\begin{figure}[H]
\centering
\begin{tikzpicture}[scale=0.75]
  \draw[-{Stealth}] (-4.6,0) -- (4.6,0) node[right] {$x$};
  \draw[-{Stealth}] (0,-4.5) -- (0,4.5) node[above] {$y$};
  % light grid
  \foreach \x in {-4,-3,-2,-1,1,2,3,4}
    \draw[gray!20] (\x,-4.3) -- (\x,4.3);
  \foreach \y in {-4,-3,-2,-1,1,2,3,4}
    \draw[gray!20] (-4.3,\y) -- (4.3,\y);
  % tick marks
  \foreach \x in {-4,-3,-2,-1,1,2,3,4}
    \draw (\x,1.5pt) -- (\x,-1.5pt) node[below, font=\scriptsize] {$\x$};
  \foreach \y in {-4,-3,-2,-1,1,2,3,4}
    \draw (1.5pt,\y) -- (-1.5pt,\y) node[left, font=\scriptsize] {$\y$};
  % clip all curves to viewport [-4,4]x[-4,4]
  \begin{scope}
    \clip (-4,-4) rectangle (4,4);
    % sin(x) — black
    \draw[black, very thick] plot[domain=-4:4, samples=120]
      (\x, {sin(\x r)});
    % T1(x) = x — red
    \draw[red, thick] plot[domain=-4:4, samples=30]
      (\x, {\x});
    % T3(x) = x - x^3/6,  Horner form: x*(1 - x^2/6) — blue
    \draw[blue, thick] plot[domain=-4:4, samples=60]
      (\x, {\x*(1 - \x*\x/6)});
    % T5(x) = x - x^3/6 + x^5/120,  Horner: x*(1 + x^2*(-1/6 + x^2/120)) — orange
    \draw[orange, thick] plot[domain=-4:4, samples=80]
      (\x, {\x*(1 + \x*\x*(-1/6 + \x*\x/120))});
  \end{scope}
  % legend
  \draw[fill=white, draw=gray!60, rounded corners=2pt]
    (2.1,2.0) rectangle (3.95,3.95);
  \draw[black, very thick] (2.25,3.7)  -- (2.75,3.7)
    node[right, font=\footnotesize] {$\sin x$};
  \draw[red, thick]        (2.25,3.2)  -- (2.75,3.2)
    node[right, font=\footnotesize] {$T_1$};
  \draw[blue, thick]       (2.25,2.7)  -- (2.75,2.7)
    node[right, font=\footnotesize] {$T_3$};
  \draw[orange, thick]     (2.25,2.2)  -- (2.75,2.2)
    node[right, font=\footnotesize] {$T_5$};
\end{tikzpicture}
\caption{Приближения на Тейлър за $\sin x$ около $x_0=0$:\; $T_1=x$;\; $T_3=x-\tfrac{x^3}{6}$;\; $T_5=x-\tfrac{x^3}{6}+\tfrac{x^5}{120}$}
\end{figure}
\end{example}

С нарастването на степента $n$ полиномът $T_n(x)$ дава все по-точно приближение на $f(x)$ в все по-голяма околност на $x_0$ — интуитивно, всяка нова производна „описва" следващо ниво на кривината на графиката.

\end{FlushLeft}

\end{document}